\section*{Chapter \ref{ch:b1:fluid-envelopes-structure} --- Structure of the Atmosphere and the Ocean}
% values: calc/ch26.py

\begin{solution}{exo:b1:fluid-envelopes-structure:1}
(a)~$H = R^*T/(M_0 g_0) = 8314.32 \times 250/(28.9644 \times 9.80665) =
\qty{7.32}{km}$. (b)~$e^{10/7.32} = 3.92$: the pressure falls to about a quarter.
(c)~$H$ is proportional to the temperature: at \qty{216.65}{K} it is
\qty{6.34}{km}, against \qty{8.43}{km} at \qty{288.15}{K}. Cold air is denser at
a given pressure, so the pressure falls faster with height.
\end{solution}

\begin{solution}{exo:b1:fluid-envelopes-structure:2}
With $\Gamma_d = \qty{9.76}{K/km}$: \qtyrange{0}{0.3}{km}, $11.0 > 9.76$,
unstable; \qtyrange{0.3}{1.2}{km}, $9.7$, neutral to within the precision of the
data; \qtyrange{1.2}{1.5}{km}, an inversion; \qtyrange{1.5}{8}{km}, stable;
above \qty{8}{km}, isothermal, stable. The inversion is the most stable, since
$\Gamma_d - \Gamma = 13.8$ is largest there.
\end{solution}

\begin{solution}{exo:b1:fluid-envelopes-structure:3}
(a)~$\sigma = 1026.95 - 1000 = \qty{26.95}{kg.m^{-3}}$. (b)~B, the lighter, lies
above A. (c)~All seawater densities begin with the same digits, 10; $\sigma$ keeps
only the digits that change from one water to another.
\end{solution}

\begin{solution}{exo:b1:fluid-envelopes-structure:4}
(a)~$S = 0.25 \times 35.40 + 0.75 \times 34.30 = 34.575$; $\theta = 0.25 \times
15.0 + 0.75 \times 4.0 = \qty{6.75}{\celsius}$. (b)~On the segment PQ, a quarter
of the way from Q to P (\cref{prop:b1:fluid-envelopes-structure:mixing}).
(c)~The mixtures of P, Q and R fill the triangle PQR. Writing
$(34.60, 8.0) = a\,\mathrm P + b\,\mathrm Q + c\,\mathrm R$ with $a + b + c = 1$
gives $a = 0.345$, $b = 0.788$, $c = -0.133$: a negative proportion of R, so the
point lies outside the triangle and no mixture of the three can produce it.
\end{solution}

\begin{solution}{exo:b1:fluid-envelopes-structure:5}
(a)~$\theta = 253.15 \times 2^{2/7} = \qty{308.6}{K}$. (b)~At \qty{1000}{hPa},
$\theta = T = \qty{298.15}{K}$. The \omterm{def:b1:fluid-envelopes-structure:potential-temperature}{potential temperature} increases upwards, from
298.2 to \qty{308.6}{K}: the layer is stable for dry air. (c)~It would arrive at
its \omterm{def:b1:fluid-envelopes-structure:potential-temperature}{potential temperature}, \qty{308.6}{K}, that is \qty{35.4}{\celsius}, ten
degrees warmer than the air already there.
\end{solution}

\begin{solution}{exo:b1:fluid-envelopes-structure:6}
(a)~$T = 288.15 - 6.5 \times 5 = \qty{255.65}{K}$ and $p = \num{101325} \times
(255.65/288.15)^{5.256} = \qty{54020}{Pa}$, \qty{540.2}{hPa}. (b)~$H =
\qty{8.43}{km}$ and $p = \num{101325}\,e^{-5/8.43} = \qty{56010}{Pa}$. The real
(standard) pressure is lower: the air there is colder than \qty{288.15}{K}, hence
denser, and the pressure falls faster with height. (c)~$p = \num{22632}\,
e^{-9/6.34} = \qty{5475}{Pa}$, the value of
\cref{tab:b1:fluid-envelopes-structure:layers}.
\end{solution}

\begin{solution}{exo:b1:fluid-envelopes-structure:7}
(a)~$\beta\,\Delta S = \alpha\,\Delta T$ gives $\Delta S = 1.67/7.61 = 0.22$.
(b)~$\Delta\rho = 1026.95 \times (-1.67\times10^{-4} \times 2 + 7.61\times10^{-4}
\times 0.3) = \qty{-0.11}{kg.m^{-3}}$: the water becomes lighter.
(c)~$\Delta\rho = 1028.11 \times (-0.53\times10^{-4} \times 2 + 7.85\times10^{-4}
\times 0.3) = \qty{+0.13}{kg.m^{-3}}$: the same changes now make it denser, because
warming barely changes the density of water near freezing.
\end{solution}

\begin{solution}{exo:b1:fluid-envelopes-structure:8}
(a)~$N^2 = (9.81/260) \times (9.76 - 4)\times10^{-3} = \qty{2.17e-4}{s^{-2}}$,
$N = \qty{0.0147}{s^{-1}}$, period $2\pi/N = \qty{426}{s}$, about seven
minutes. (b)~$N^2 = (9.81/270) \times 14.76\times10^{-3} = \qty{5.36e-4}{s^{-2}}$,
$N = \qty{0.0232}{s^{-1}}$, period \qty{271}{s}. (c)~$N$ tends to zero and the
period grows without bound: the restoring force vanishes in a neutral layer.
\end{solution}

\begin{solution}{exo:b1:fluid-envelopes-structure:9}
(a)~It is denser, $\sigma_\theta = 27.20$ against 26.30: the AAIW is fresher
but much colder. (b)~$\Delta\sigma = 1026.95 \times [-1.67\times10^{-4} \times
(4.31 - 15.43) + 7.61\times10^{-4} \times (34.30 - 35.54)] = 1.907 - 0.969 =
0.94$, against $0.90$ exactly. (c)~Temperature: the cooling adds
\qty{1.91}{kg.m^{-3}}, the freshening removes \qty{0.97}{kg.m^{-3}}.
\end{solution}

\begin{solution}{exo:b1:fluid-envelopes-structure:10}
(a)~By \cref{thm:b1:fluid-envelopes-structure:stability}, $(g/\theta)\,\mathrm
d\theta/\mathrm dz = (g/T)(\Gamma_d - \Gamma)$; $\mathrm d\theta/\mathrm dz = 0$
gives $\Gamma = \Gamma_d$. (b)~$T = T_s - \Gamma_d z$, a layer of constant gradient
$L = \Gamma_d$: by \cref{prop:b1:fluid-envelopes-structure:layer-pressure}, $p =
p_s (T/T_s)^{gM/(R\Gamma_d)}$, and $gM/(R\Gamma_d) = gMc_p/(Rg) = c_pM/R$
($= 3.5$ for dry air). (c)~$p$ vanishes where $T = 0$, at $z = T_s/\Gamma_d =
300/9.76 = \qty{30.7}{km}$. Real layers of uniform \omterm{def:b1:fluid-envelopes-structure:potential-temperature}{potential temperature} are a
few kilometres deep at most: they are capped by stable air, as the troposphere is
by the stratosphere.
\end{solution}

\begin{solution}{exo:b1:fluid-envelopes-structure:11}
(a)~$S = (33.787 + 36.050)/2 = 34.919$, $\theta = \qty{8.0}{\celsius}$.
(b)~With a linear equation of state, $\sigma$ is an affine function of $S$ and
$\theta$, so the $\sigma$ of the mixture is the mean of the two, 27.00.
(c)~$\alpha$ is \qty{0.75e-4}{K^{-1}} at \qty{2}{\celsius} and
\qty{2.07e-4}{K^{-1}} at \qty{14}{\celsius}, while $\beta$ changes little
(\num{7.80e-4} and \num{7.53e-4}). The density falls ever faster as the
temperature rises: it is a concave function of temperature. Along the mixing
line, the mixture is at the mean temperature, where the density lies above the
chord joining the two parents, that is above 27.00. (d)~The excess comes from the
curvature, a second-order term: it grows as the square of the temperature
difference of the parents.
\end{solution}

\begin{solution}{exo:b1:fluid-envelopes-structure:12}
(a)~The pressure difference is $5604 - 5500 = \qty{104}{dbar}$: the parcel cools
by $0.129 \times 0.104 = \qty{0.013}{\celsius}$, to \qty{0.718}{\celsius}.
(b)~The water around it at \qty{5400}{m} is at \qty{0.730}{\celsius}, warmer: the
parcel, of nearly the same salinity, is denser than its surroundings and sinks
back. The layer is stable, although the measured temperature increases
downwards; its \omterm{def:b1:fluid-envelopes-structure:potential-temperature}{potential temperature} decreases downwards. (c)~$\theta \approx
0.7313 - 0.086 \times 5.604 = \qty{0.249}{\celsius}$, within a thousandth of a
degree of the exact value.
\end{solution}

\begin{solution}{pb:b1:fluid-envelopes-structure:1}
\begin{enumerate}
  \item The temperature is nearly uniform to \qty{50}{m}, drops fast to
        \qty{300}{m}, more slowly to \qty{1000}{m}, and is nearly uniform below.
        The salinity decreases from 35.80 at the surface to a minimum of 34.32 at
        \qty{800}{m}, then increases to 34.90 at \qty{3000}{m}.
  \item $22.40 - 0.5 = \qty{21.90}{\celsius}$ is reached between \qty{50}{m}
        (22.10) and \qty{75}{m} (19.20): $50 + 25 \times 0.20/2.90 = \qty{52}{m}$.
  \item $(22.10 - 13.00)/250 = \qty{0.0364}{\celsius/m}$, \qty{3.6}{\celsius}
        per \qty{100}{m}. The \omterm{def:b1:fluid-envelopes-structure:ocean-layers}{thermocline} extends from the base of the \omterm{def:b1:fluid-envelopes-structure:ocean-layers}{mixed
        layer}, near \qty{50}{m}, to about \qty{800}{m}; it is sharpest between
        \qty{50}{m} and \qty{75}{m} (\qty{11.6}{\celsius} per \qty{100}{m}).
  \item $\sigma = 1026.95 \times [1 - 1.67\times10^{-4} \times 12.4 +
        7.61\times10^{-4} \times 0.80] - 1000 = 25.45$, \qty{0.74}{kg.m^{-3}} too
        high. The surface water is far from the reference state, and $\alpha$
        is much larger at \qty{22}{\celsius} than at \qty{10}{\celsius}: the
        linear form underestimates the lightening by the warmth.
  \item The salinity minimum at \qty{800}{m}, 34.32 at \qty{4.8}{\celsius}, is
        Antarctic Intermediate Water, formed under the rainy westerlies of the
        Southern Ocean. The maximum at \qty{3000}{m}, 34.90 at
        \qty{2.5}{\celsius}, is North Atlantic Deep Water, formed in the
        Greenland and Labrador Seas.
  \item $(2.50 - 2.26)/3.042 = \qty{0.079}{\celsius}$ per \qty{1000}{dbar}.
  \item $\sigma_\theta$ increases with depth: the column is stable. The
        \omterm{def:b1:fluid-envelopes-structure:ocean-layers}{pycnocline} is where $\sigma_\theta$ increases fastest, in the upper few
        hundred metres, with the \omterm{def:b1:fluid-envelopes-structure:ocean-layers}{thermocline}: \num{0.82} between
        \qty{200}{m} and \qty{800}{m}, only \num{0.70} over the
        \qty{2200}{m} below.
  \item $S = (35.45 + 34.32)/2 = 34.885$; $\theta = (14.90 + 4.80)/2 =
        \qty{9.85}{\celsius}$.
  \item C: 26.461; I: 27.310; mixture: 26.886, exactly the mean of the two
        (26.886), since the linear form is affine.
  \item The mean of the parents is $(26.342 + 27.159)/2 = 26.751$; the mixture,
        at 26.888, is \qty{0.14}{kg.m^{-3}} denser. This is cabbeling.
  \item With the linear form, the cast gives 26.61 at \qty{300}{m}, 26.88 at
        \qty{500}{m} and 27.31 at \qty{800}{m}: the mixture's density lies near
        \qty{500}{m}. The water there, $(34.82, \qty{9.6}{\celsius})$, is very
        close to the mixture $(34.885, \qty{9.85}{\celsius})$: the cast between
        the Central Water and the AAIW is itself close to a mixing line between
        them.
  \item $S = 0.25 \times 35.45 + 0.75 \times 34.32 = 34.60$, $\theta = 0.25
        \times 14.90 + 0.75 \times 4.80 = \qty{7.33}{\celsius}$. Interpolating
        the cast between \qty{500}{m} and \qty{800}{m} gives
        $(34.57, \qty{7.2}{\celsius})$ near \qty{650}{m}: nearly on it.
  \item In \si{K/km}: A--B 10.0; B--C 9.7; C--D $-5.5$; D--E 6.29; E--F 6.50;
        F--G 6.57; G--H 0.33; H--I $-0.50$; I--J $-0.83$.
  \item A--B unstable; B--C neutral; C--D an inversion; D--E, E--F, F--G stable;
        G--H nearly isothermal, stable; H--I and I--J inversions. The tropopause
        is at G, \qty{12.5}{km} and \qty{-59.0}{\celsius}, where the temperature
        stops decreasing.
  \item The mean temperature is \qty{214.65}{K}: $H = 287 \times 214.65/9.81 =
        \qty{6.28}{km}$, and $147\,e^{-4/6.28} = \qty{77.7}{hPa}$, consistent
        with the 78 measured.
  \item The air at A, \qty{22.0}{\celsius}, is cooler than the sea: \omterm{def:b1:heat-budget:heat-flow}{heat flows}
        from the sea to the air, which is heated from below. The layer A--B, with
        a lapse rate above $\Gamma_d$, is the result.
  \item At B: $22.0 - 9.76 \times 0.5 = \qty{17.1}{\celsius}$, against 17.0; at
        C: \qty{7.4}{\celsius}, against 7.3; at D: \qty{3.5}{\celsius}, against
        9.5. The parcel is slightly warmer than the air up to C, and much colder
        at D.
  \item At the base of the inversion, near \qty{1.5}{km}: entering the inversion,
        the parcel becomes colder, hence denser, than the air around it, which
        stops its rise.
  \item Parcels released by the unstable surface layer rise through it and mix it
        until its \omterm{def:b1:fluid-envelopes-structure:potential-temperature}{potential temperature} is uniform, which makes its lapse rate
        $\Gamma_d$; the inversion above stops the mixing.
  \item In K: A 294.15; B 293.99; C 294.08; D 300.51; E 312.65; F 329.29; G
        346.28; H 369.50; I 446.99; J 600.91.
  \item $\theta$ decreases from A to B (unstable), is nearly constant from B to
        C (neutral) and increases above (stable), fastest from C to D and above
        G: the same classification as in question~14.
  \item E--F: $(329.29 - 312.65)/4 = \qty{4.2}{K/km}$; H--I: $(446.99 -
        369.50)/4 = \qty{19.4}{K/km}$, nearly five times more.
  \item With the mean $\theta = \qty{408.2}{K}$: $N^2 = (9.81/408.2) \times
        0.01937 = \qty{4.66e-4}{s^{-2}}$, $N = \qty{0.0216}{s^{-1}}$, a period
        of \qty{291}{s}, about five minutes.
  \item $N^2 = (9.81/1027) \times (26.35 - 24.79)/150 = \qty{9.9e-5}{s^{-2}}$,
        $N = \qty{0.0100}{s^{-1}}$, a period of \qty{630}{s}. The lower
        stratosphere is the more stable, with $N^2$ nearly five times larger
        than in the \omterm{def:b1:fluid-envelopes-structure:ocean-layers}{thermocline}.
  \item At G, $\theta = 214.15 \times (1000/186)^{2/7} = \qty{346}{K}$: the
        \omterm{def:b1:fluid-envelopes-structure:potential-temperature}{potential temperature} at the tropopause is about \qty{346}{K}.
\end{enumerate}
\end{solution}
